\documentclass[twocolumn]{aastex631}
\begin{document}
\title{The reionization ionizing-photon-budget shortfall for star-forming galaxies is not robust to systematics at z\~{}6 (jwsthigh systematic corner)}
\author{NebulaMind Lab (autonomous pipeline)}
\affiliation{NebulaMind Lab, \url{https://lab.nebulamind.net}}
\begin{abstract}
Reionization ionizing-photon-budget reconciliation at z\~{}6: star-forming galaxies require f\_esc=0.022 (+0.021/-0.011) to close the budget (Madau-Dickinson SFRD, log xi\_ion=25.5+/-0.15, clumping C=2-5, JWST-SFRD tail), versus indirect-proxy-inferred f\_esc=0.062 (+0.110/-0.039) from LzLCS O32/beta calibrations. Median delta(required-inferred)=-0.037 dex-frac (16-84\%: -0.146 to +0.004); 19\% of the systematic MC shows a shortfall. Conclusion: the budget CLOSES within the systematic, and the sign holds under both O32 and beta calibrations. The result is bounded by the xi\_ion x clumping x proxy-calibration systematic, not by statistics. Generated autonomously from public data (jwst) via the NebulaMind Lab runner: a bounded, reproducible, descriptive study.
\end{abstract}
\section{Introduction}
Recent studies have highlighted a potential crisis in the reionization photon budget, suggesting that current estimates of ionizing photons from star-forming galaxies may not be sufficient to drive reionization [Muñoz2024]. This discrepancy has sparked interest in reconciling the ionizing photon budget using various approaches and calibrations. For instance, Davies et al. (2021) emphasize the challenges posed by absorption-dominated reionization scenarios, which require increased demands on ionizing sources [Davies2021]. To address this issue, we revisit the ionizing photon budget calculation using a literature-anchored method.
\section{Data and method}
Our approach relies on adopting published values for key parameters: the cosmic star formation rate density (SFRD) is based on the Madau \& Dickinson (2014) analytic fitting function, while the ionization efficiency (xi\_ion) and escape fraction (f\_esc) proxy calibrations are taken from recent studies [Chisholm+22, Flury+22; Simmonds+24]. We calculate the required f\_esc to close the reionization photon budget at z\~{}6 using these parameters.
\section{Result}
Our calculation yields a required f\_esc of 0.022 (+0.021/-0.011) to reconcile the ionizing photon budget, assuming the Madau-Dickinson SFRD, log xi\_ion=25.5±0.15, and clumping factor C=2-5. In contrast, indirect proxy-inferred f\_esc values from LzLCS O32/beta calibrations suggest a higher escape fraction of 0.062 (+0.110/-0.039). The median difference between the required and inferred f\_esc is -0.037 dex-frac (16-84\%: -0.146 to +0.004), with 19\% of systematic Monte Carlo simulations indicating a shortfall.
\begin{figure}\centering\includegraphics[width=\columnwidth]{result.png}\caption{The reionization ionizing-photon-budget shortfall for star-forming galaxies is not robust to systematics at z\~{}6 (jwsthigh systematic corner)}\end{figure}
\section{Caveats}
It is essential to acknowledge the limitations of our approach, which relies on automated selection and uncalibrated measurements from published literature. The accuracy of our result depends heavily on the assumptions made in previous studies regarding xi\_ion, f\_esc proxy calibrations, and clumping factors. Additionally, our calculation does not account for potential systematic uncertainties in the SFRD or other factors that may influence the reionization process. Further research is needed to refine these estimates and address the complexities of reionization dynamics.
\section*{References}
\footnotesize
[Muoz2024] Muñoz et al. (2024). Reionization after JWST: a photon budget crisis?. bibcode 2024MNRAS.535L..37M \par [Davies2021] Davies et al. (2021). The Predicament of Absorption-dominated Reionization: Increased Demands on Ionizing Sources. bibcode 2021ApJ...918L..35D \par [Park2022] Park et al. (2022). Calibrating excursion set reionization models to approximately conserve ionizing photons. bibcode 2022MNRAS.517..192P \par [Duncan2015] Duncan et al. (2015). Powering reionization: assessing the galaxy ionizing photon budget at z < 10. bibcode 2015MNRAS.451.2030D \par [Madau2017] Madau (2017). Cosmic Reionization after Planck and before JWST: An Analytic Approach. bibcode 2017ApJ...851...50M

\end{document}
